\section*{Preface}
\addcontentsline{toc}{section}{Preface}

This document was written for one reader: someone sitting down to reproduce \textcite{fountas2025emllm} and wanting a single place that explains both the paper's machinery and the code that implements it. It assumes familiarity with transformer attention, with the basics of long-context inference, and with the conventions of a Hugging Face training stack. It does not assume prior exposure to event-cognition theory, graph-based clustering metrics, or the particular abstractions used inside \file{em\_llm/}.

The chapters are ordered to be read straight through, but each is intended to stand on its own when used as a reference. Part I covers the paper. Chapter~1 places the work in context and gives the headline results. Chapter~2 builds the cognitive-science background that makes the rest of the paper legible. Chapter~3 is the longest of Part I and steps through the three core innovations (surprise-based segmentation, graph-theoretic boundary refinement, two-stage retrieval) with worked examples. Chapters~4 and~5 cover architecture, hyperparameters, and the experimental setup. Chapter~6 closes Part I with limitations and open questions.

Part II is the code companion. Chapter~7 maps the repository so the reader knows what is vendored, what is reproduction-specific, and how OmegaConf layers configuration at runtime. Chapter~8 is the function-by-function walk through \file{em\_llm/}, which is where the paper's algorithm actually lives. Chapter~9 covers the benchmark harness, and Chapter~10 covers the reproduction scripts in \file{reproduction/}. Two appendices follow: a glossary for the cognitive-science and graph-theory vocabulary, and a notation table.

A note on the relationship between this document and the paper. Anywhere the explanation here improves on or diverges from the paper's phrasing, it is in service of the reader, not as a correction. Where the paper introduces a symbol, a section, a figure, or an algorithm, the companion references it explicitly so the reader can flip to the original. Where the code uses a name that differs from the paper, both names appear together the first time the concept is introduced.
